\uuid{P7nd}
\exo7id{4151}
\titre{exo7 4151}
\auteur{quercia}
\organisation{exo7}
\datecreate{2010-03-11}
\isIndication{false}
\isCorrection{true}
\chapitre{Fonction de plusieurs variables}
\sousChapitre{Dérivée partielle}
\cpge{mpsi-mp2i.fonctions-deux-variables.b, pcsi.fonctions-deux-variables.b, ptsi.fonctions-deux-variables.b, mpsi-mp2i.derivabilite.d, pcsi.derivabilite.d, ptsi.derivabilite.d, bcpst1.fonctions-deux-variables.c, bcpst1.fonctions-deux-variables.d, mp-mpi.calcul-differentiel-optimisation.a}
\module{Analyse}
\niveau{L2}
\difficulte{}

\contenu{
\question{
Soient ${f,g} : {\R^2} \to \R$ de classe $\mathcal{C}^n$. Calculer
$\frac{\partial^n(fg)}{\partial x^k\partial y^{n-k}}$ en fonction des dérivées
de $f$ et $g$.
}
\reponse{
$\frac{\partial^n(fg)}{\partial x^k\partial y^{n-k}} =
\sum_{i=0}^k\sum_{j=0}^{n-k} C_k^iC_{n-k}^j
\frac{\partial^{i+j}f}{\partial x^i\partial y^j}
\frac{\partial^{n-i-j}g}{\partial x^{k-i}\partial y^{n-k-j}}$.
}
}
